\uuid{8wPW}
\exo7id{2469}
\titre{exo7 2469}
\auteur{matexo1}
\organisation{exo7}
\datecreate{2002-02-01}
\isIndication{false}
\isCorrection{false}
\chapitre{Réduction d'endomorphisme, polynôme annulateur}
\sousChapitre{Diagonalisation}
\cpge{mpsi-mp2i.applications-lineaires.b, pcsi.applications-lineaires.b, ptsi.applications-lineaires.b, mpsi-mp2i.changements-bases.c, mpsi-mp2i.determinants.f, mp-mpi.reduction-endomorphismes.b, pc.reduction.a, psi.reduction-endomorphismes.a, pt.reduction-endomorphismes.a, mp-mpi.reduction-endomorphismes.c, pc.reduction.b, psi.reduction-endomorphismes.b, pt.reduction-endomorphismes.b, mp-mpi.reduction-endomorphismes.d, pc.reduction.c, psi.reduction-endomorphismes.c, pt.reduction-endomorphismes.c, bcpst2.valeurs-propres-vecteurs-propres.a, bcpst2.valeurs-propres-vecteurs-propres.b}
\module{Algèbre}
\niveau{L2}
\difficulte{}

\contenu{
\question{
Soit $a,b,c$ trois r\'eels et $u$ l'endomorphisme de $\R^3$
ayant pour matrice
$$M = \left( \begin{array}{ccc} 1&a&b\\ 0&1&c \\ 0&0&-1 \end{array} \right).$$
Discuter la possibilit\'e de le diagonaliser selon les valeurs de
$a,b,c$.
}
}
